%%
%% winter-lecture27.tex
%% 
%% Made by Alex Nelson <pqnelson@gmail.com>
%% Login   <alex@lisp>
%% 
%% Started on  2026-03-07T14:44:19-0800
%% Last update 2026-03-07T14:44:19-0800
%% 

\lecture{}

\begin{node}
We will try to get to proving the Bott periodicity theorem. There are
two proofs:
\begin{enumerate}
\item Bott's approach using infinite-dimensional Morse theory, and
\item Hatcher's approach using K-Theory and clutching construction.
\end{enumerate}
We will probably flesh out Hatcher's proof.
\end{node}

\begin{remark}
In a principal $G$-bundle, we think of the fiber $G$ as a $G$-torsor
(i.e., we forgot about the identity element).
\end{remark}

\begin{proposition}
Given $A\stackrel{i}{\into}X\xrightarrow{q}X/A$ when $X$ is comapct
Hausdorff and $A\subset X$ is closed, we have
\begin{equation}
\widetilde{K}(X/A)\xrightarrow{q^{*}}\widetilde{K}(X)\xrightarrow{i^{*}}\widetilde{K}(A)
\end{equation}
is exact.
\end{proposition}

\begin{proof}
We want to show $\ker(i^{*})=\Im(q^{*})$.

\textsc{Claim 1}: $\ker(i^{*})\supset\Im(q^{*})$. That is to say,
$i^{*}\circ q^{*}=0$. We start with a bundle $E\to X/A$, then the
pullback bundle $q^{*}E\to X$ has the property that $(q^{*}E)|_{A}$ is
trivial over $A$. Then $i^{*}(q^{*}E)=q^{*}E|_{A}$ which is a trivial
bundle. Then $[i^{*}(q^{*}E),\CC^{N}]=0\in\widetilde{K}(A)$ where $N=\dim(E)$.

\textsc{Claim 2}: $\ker(i^{*})\subset\Im(q^{*})$. Start with $E\to X$
which is in $\ker(i^{*})$. Then $i^{*}E=E|_{A}$ is stably trivial,
i.e., there exists some $n\in\NN$ such that
\begin{equation}
E|_{A}\oplus\TrivialBundle{\CC}^{n}|_{A}\iso\TrivialBundle{\CC}^{N}|_{A},
\end{equation}
so replacing $E$ by $E\oplus\TrivialBundle{\CC}^{n}|_{E}$ we can
assume without loss of generality $E|_{A}$ is trivial in the sense
that there exists a local trivialization of $E|_{A}$,
\begin{equation}
h\colon A\times\CC^{n}\to E|_{A}.
\end{equation}
Define $E/h$ to be the quotient space of $E$ by the relation
$\forall x,y\in A\forall v\in\CC^{n}\ldotp h(x,v)\sim h(y,v)$,
that is to say we identify all the fibers. Then there is a projection
\begin{equation}
p\colon E/h\to X/A
\end{equation}
such that the fiber over each point is $\CC^{n}$ and $q^{*}(E/h)\iso E$.
We just need to show $p\colon E/h\to X/A$ is a vector bundle. We need
to prove local triviality, the tricky bit is the point $A/A$. In
particular, we want local trivial charts containing the point $A/A\in X/A$.

\textsc{Sketchy proof sketch}. We know $E|_{A}$ is trivial. Then there
exists nowhere dependent sections $s_{1}$, \dots, $s_{n}$ of
$E|_{A}$. We want to find an open neighborhood $U\supset A$. Since $A$
is closed, we can extend $s_{1}$, \dots, $s_{n}$ to sections of $E$
over $X$---you need to use compactness arguments. Note that being
linearly independent is an \emph{open condition} (i.e., the sections
$s_{1}$, \dots, $s_{n}$ gives us a matrix
$A=(s_{1}|\dots|s_{n})\in\GL_{n}(\CC)$ whose column vectors are the
sections, and then linear independence is equivalent to $\det(A)\neq0$).
Since the $s_{1}$, \dots, $s_{n}$ are linearly independent over $A$,
then there exists some open neighborhod $U\supset A$ such that these
sections are still linearly independent over $U$. This means $E|_{U}$
is trivial as well, and this gives the trivial chart in the quotient
we wanted.
\end{proof}

\begin{node}
Exactness in the middle gives us everything we want. Consider the
Puppe sequence of cofibrations---recall, given a CW complex $i\colon A\into X$
we have the Puppe sequence look like
\begin{equation}
A\stackrel{i}{\into}X\xrightarrow{q}\to X/A\to\Sigma A\to\Sigma X\to\Sigma(X/A)\to\Sigma^{2}(A)\to\dots
\end{equation}
such that every 3 adjacent sequential terms looks like $U\into V\to V/U$
for some CW pair $(U,V)$.

Therefore, applying this to $[-,BU]$ gives us a long exact sequence of
the form
\begin{equation}
\dots\to\widetilde{K}(\Sigma X)\to\widetilde{K}(\Sigma A)\to\widetilde{K}(X/A)\to\widetilde{K}(X)\to\widetilde{K}(A)
\end{equation}
which looks like the long exact sequence for a cohomology theory.

Recall for the usual cohomology theory,
\begin{equation}
\widetilde{H}^{n}(\Sigma X)\iso\widetilde{H}^{n-1}(X).
\end{equation}
We want to mimic this behaviour. If we denote the usual $K$-theory as
\begin{equation}
\widetilde{K}^{0}(X):=\widetilde{K}(X),
\end{equation}
then $\widetilde{K}(\Sigma X)$ should be thought of as $\widetilde{K}^{-1}(X)$.
\end{node}

\begin{definition}
Let $X$ be a compact Hausdorff space.
For $n\in\NN_{0}$, the \define{$(-n)^{\text{th}}$ Reduced K-Theory} of $X$
is $\widetilde{K}^{-n}(X):=\widetilde{K}(\Sigma^{n}X)$.
\end{definition}

\begin{question}
What happens for the positive degree? That is to say, what are the
$\widetilde{K}^{n}(X)$ for $n\in\NN$?
\end{question}

\begin{theorem}[Bott periodicity for topological complex K-Theory]
We have $\widetilde{K}(X)\iso\widetilde{K}(\Sigma^{2}X)$ with an
explicit map called the \define{Bott Map}
\begin{equation}
\begin{split}
\beta\colon&\widetilde{K}(X)\to\widetilde{K}(\Sigma^{2}X)\iso\widetilde{K}(\sphere{2}\smashtimes X)\\
&E\mapsto(H-1)\boxtimes E
\end{split}
\end{equation}
\end{theorem}

\begin{node}
From this perspective, there are only 2 nontrivial pieces of
information: $\widetilde{K}^{0}(X)$ and $\widetilde{K}^{1}(X)$. Then
we can define
\begin{equation}
\widetilde{K}^{n}(X)=\begin{cases}\widetilde{K}^{0}(X) & \mbox{if $n$
  is even}\\
\widetilde{K}^{1}(X) & \mbox{if $n$ is odd}
\end{cases}
\end{equation}
for any $n\in\ZZ$.
\end{node}

\subsection{Examples of relative exact sequences}

\begin{node}
If $A\into X$ where $A$ is a sub-CW complex of $X$, and if $A$ is
contractible, then
\begin{equation}
\widetilde{K}(\Sigma A)=0\to\widetilde{K}(X/A)\to\widetilde{K}(X)\to\widetilde{K}(A)\iso\widetilde{K}(\point{*})=0,
\end{equation}
so $q^{*}\colon\widetilde{K}(X/A)\to\widetilde{K}(X)$ is an isomorphism.
\end{node}

\begin{proof}
Contractible $A$ has $\widetilde{K}(A)\iso\widetilde{K}(\point{*})=0=0$.
Since $A$ is contractible, $\Sigma A$ is contractible, so $\widetilde{K}(\Sigma A)=0$.
Hence the claim.
\end{proof}

\begin{remark}
This result still holds in the more general setting where $A\into X$
is a closed subspace and $A$ is contractible, then we always have this
isomorphism on $\widetilde{K}$ even though $X/A$ is not necessarily
homotopy equivalent to $X$.
\end{remark}

\begin{example}
The wedge product $X\wedgetimes Y$ satisfies
\begin{equation}
\widetilde{K}(X\wedgetimes Y)\iso\widetilde{K}(X)\oplus\widetilde{K}(Y).
\end{equation}
\end{example}

\begin{definition}
The smash product $X\smashtimes Y$ and external product
$\boxtimes$. Given pointed spaces $(X,x_{0})$ and $(Y,y_{0})$ we can
form the product $(X\times Y,x_{0}\times y_{0})$, but this is not the
product object in the category of pointed topological spaces $\TOP_{*}$.
This product object is described by the more natural notion of the
\define{Smash Product}
\begin{equation}
X\smashtimes Y:=\frac{X\times Y}{(\{x_{0}\}\times Y)\cup(X\times\{y_{0}\})}.
\end{equation}
\end{definition}

\begin{example}
$\sphere{1}\smashtimes\sphere{1}\iso\sphere{2}$
\end{example}

\begin{example}
$\sphere{i}\smashtimes\sphere{j}\iso\sphere{i+j}$
\end{example}

\begin{example}
$\sphere{i}\smashtimes X\iso\Sigma X$ is the reduced suspension of $X$.
The intuition is that we think of $\sphere{1}$ as the unit interval
with the endpoints identified $\sphere{1}\iso I/(0\sim 1)$. 
\end{example}